\begin{abstract}

Short-text classification in Portuguese — such as product descriptions from electronic invoices, with abbreviated, noisy vocabulary and hundreds of categories — depends on large volumes of labels whose acquisition is the dominant cost of these systems. This thesis proposes and evaluates FALCO (Active Learning Framework with LLMs for Short Text), which integrates three cost-reduction fronts: informed composition of the initial training set, uncertainty-based active selection of instances, and partial replacement of the human annotator by oracles based on large language models (LLMs), organized in phases of increasing cost. The investigation uses an audited public dataset of 250,365 product descriptions across 621 categories (250,221 descriptions in the corrected version used in the experiments) and follows a pre-registered protocol with immutable partitions, paired analysis (McNemar, Wilcoxon), and traceable artifacts for every reported number. Four main results: (i) sensitivity to the composition of the initial set is highest in the low-label regime (a 6.4 percentage-point accuracy range at $|L_0|=100$), and the proposed deterministic heuristic, DRI-SL — semantic density with lexical diversity, using no labels at all — outperforms the best individual found by a supervised genetic algorithm at every size from 100 to 5,000, a conservative comparison after the circularity correction that reduced the evolutionary envelope by up to 6.3 percentage points; (ii) zero-shot LLM oracles reach 77--83\% accuracy in the closed space of 621 categories at costs of US\$~0.035 to US\$~0.92 per thousand labels, with Macro F1 above that of a light supervised classifier trained on 250 thousand labels, and nearly half of their errors concentrate on catalog conventions (sibling categories and umbrella classes); (iii) the measurement instrument matters: without output restriction to the class space, 6.8\% of responses become spurious errors; boundary rules in the prompt improve the weak model on the production distribution ($+4.6$ p.p., $p=0.012$) at the cost of severe degradation on rare classes ($p<0.001$); and even the serving provider is part of the instrument — the same free model diverges significantly from itself across two services ($p<0.001$), although, properly served, it ties with paid oracles on the production sample; (iv) uncertainty strategies outperform random selection at the maximum significance achievable with 8 seeds ($p=0.0078$), recovering 78\% of the supervised ceiling's Macro F1 with 15\% of the labels, and this advantage survives intact under uniform oracle noise up to $\varepsilon=0.4$ — covering the error range observed in the real LLMs. Since the best oracle fell short of the pre-registered 85\% accuracy criterion, the methodological gate defined the final configuration (an economical oracle in the initial phase and the most accurate oracle in the advanced phase) and made learning with noisy labels the central scenario. Finally, (v) the framework was executed end to end with a free LLM oracle at zero monetary cost, stopping by stagnation at 32--40\% of budget, and validation with the strong classifier (BERTimbau) trained outside the loop answered the pre-registered central hypothesis — $\ge 95\%$ of full-pool-supervision Macro F1 with $\le 30\%$ of the labels —, which was \textbf{refuted at the pre-registered 30\% budget but sustained from $\approx$50\% of the pool}: the paired decomposition exonerates oracle noise (which, being structured, even benefits rare classes) and selection, attributing the loss to the stopping criterion inherited from the light classifier; the budget sweep — a \textit{post hoc} finding, distinct from the pre-registered criterion — locates the floors — 40\% of the labels for accuracy and 50\% for Macro F1 — and shows that, with 70\% of actively selected labels, the strong model surpasses full-pool supervision. Population-scale experiments further show that self-evaluation on actively collected data is biased in a metric-dependent direction — internal accuracy underestimates and internal Macro F1 overestimates real performance —, grounding the reserved evaluation set and release criterion the framework adopts. The first-order limitations are declared: the results come from a single dataset, from the author's own study domain; the decisive arm of the strong-classifier validation ran on a single seed; and the cost analysis rests on ratios between oracles — such as the 26-fold difference under a statistical accuracy tie — rather than on absolute prices, which date quickly. The thesis also delivers the LCE metric for learning-curve comparison, the open-source \texttt{activelearning} library, and the FlowBuilder tool, with sanitized data ingestion, parameterized execution, and learning curves, which make the protocol reusable in other domains.

\end{abstract}
